\section[Rédaction Bourbaki \redac{469} : le foncteur constant crée les limites]{Rédaction Bourbaki \redac{469} : le foncteur constant crée les limites\protect\piste\protect\verifie{Redaction469}}\label{sec:bourbaki-469}

La rédaction \redac{469}, \q{Catégories abéliennes}, des archives Bourbaki (Archives Henri Poincaré, \ahp{1042}) est datée de mai--juin 1966 par sa notice, qui la donne à Pierre Deligne ; elle compte 54 pages. Son théorème~4 (p.~23) montre qu'une catégorie de foncteurs $\mathcal C \to \mathcal A$ hérite des propriétés de $\mathcal A$ : limites, colimites, objet nul, structure abélienne. Les errata de la page~52 reprennent la réciproque par un lemme sur le foncteur des diagrammes constants. Sa transcription reste locale pour l'instant ; on n'en cite ici aucun passage.

\emph{La réciproque et son piège.} On voudrait conclure dans l'autre sens : si $\mathcal C \to \mathcal A$ a des limites, $\mathcal A$ en a aussi, pourvu que $\mathcal C$ ne soit pas vide. L'idée naturelle est de plonger $\mathcal A$ dans la catégorie de foncteurs par les foncteurs constants, $\mathrm{const} : \mathcal A \to \mathcal A^{\mathcal C}$, de prendre la limite $L$ du diagramme constant $\mathrm{const} \circ F$ pour un diagramme $F : J \to \mathcal A$, puis de l'évaluer en un objet $d$ de $\mathcal C$. Mais l'évaluation ne préserve les limites que si celles de la catégorie de foncteurs se calculent point par point, ce qui suppose que $\mathcal A$ en ait. L'argument tourne en rond, et il faut une autre preuve que $L(d)$ est une limite de $F$.

\begin{encadre}
\begin{enonce}[Bourbaki, rédaction \redac{469}, lemme des errata, p.~52]
Soit $\mathcal C$ une catégorie non vide. Le foncteur $\mathrm{const} : \mathcal A \to \mathcal A^{\mathcal C}$ crée les limites et les colimites : si $L$ est une limite de $\mathrm{const} \circ F$ dans $\mathcal A^{\mathcal C}$, alors, pour tout objet $d$ de $\mathcal C$, $L(d)$ est une limite de $F$ dans $\mathcal A$, et $L$ envoie toute flèche de $\mathcal C$ sur un isomorphisme. En conséquence, $\mathcal A$ a les limites (finies, ou d'une forme donnée) et les colimites que $\mathcal A^{\mathcal C}$ possède, un objet nul si $\mathcal A^{\mathcal C}$ en a un, et, si $\mathcal A$ est préadditive et $\mathcal A^{\mathcal C}$ abélienne, $\mathcal A$ est abélienne.
\end{enonce}
\end{encadre}

\emph{L'idée.} La limite $L$ est un foncteur, et chacune de ses valeurs $L(d)$ porte un cône sur $F$. Ce cône, rendu constant, donne une flèche $u_d$ du foncteur constant de valeur $L(d)$ vers $L$. On recolle ensuite les composantes diagonales $(u_d)_d$ en un endomorphisme de $L$, compatible aux projections ; par l'unicité dans la propriété universelle de $L$, c'est l'identité. Cette identité sur la diagonale suffit à donner, en chaque $d$, l'unicité qui manquait.

\begin{figure}[htbp]
\centering
\begin{tikzcd}[column sep=large]
\mathrm{const}\, L(d) \arrow[r, "u_d"] \arrow[rr, bend right=18, "\mathrm{const}\,(\pi_j)_d"'] & L \arrow[r, "\pi_j"] & \mathrm{const}\, F(j)
\end{tikzcd}
\caption{La flèche $u_d$ de la preuve, dans la catégorie de foncteurs $\mathcal A^{\mathcal C}$. Les composantes $(\pi_j)_d : L(d) \to F(j)$ forment un cône sur $F$ ; rendu constant, il se factorise par la limite $L$ de $\mathrm{const} \circ F$. Les composantes diagonales $(u_d)_d : L(d) \to L(d)$ forment un endomorphisme de $L$ compatible aux projections $\pi_j$, qui est donc l'identité.}
\label{fig:bourbaki-469}
\end{figure}

\begin{proof}
Notons $\pi_j : L \to \mathrm{const}\, F(j)$ les projections de la limite. En un objet $d$, les composantes $(\pi_j)_d : L(d) \to F(j)$ forment un cône sur $F$. Le cône constant correspondant se factorise par $L$ : il existe $u_d : \mathrm{const}\, L(d) \to L$ avec $\pi_j \circ u_d = \mathrm{const}\,(\pi_j)_d$ (figure~\ref{fig:bourbaki-469}). Pour une flèche $f : d \to d'$, les deux flèches $u_{d'} \circ \mathrm{const}\, L(f)$ et $u_d$ ont les mêmes composés avec les projections, donc sont égales. Cette relation dit que les composantes $(u_d)_d : L(d) \to L(d)$ forment une transformation naturelle $w : L \to L$. Elle vérifie $\pi_j \circ w = \pi_j$ pour tout $j$, donc $w = \mathrm{id}_L$ : chaque $(u_d)_d$ est l'identité de $L(d)$.

Montrons que $L(d)$ est une limite de $F$. Un cône $s$ de sommet $X$ sur $F$ donne, rendu constant, une flèche $\mathrm{const}\, X \to L$, dont la composante en $d$ est une flèche $X \to L(d)$ compatible au cône. Pour l'unicité, soient $m, m' : X \to L(d)$ deux flèches compatibles. Les flèches $u_d \circ \mathrm{const}\, m$ et $u_d \circ \mathrm{const}\, m'$ de $\mathrm{const}\, X$ vers $L$ ont les mêmes composés avec les projections, donc sont égales ; leurs composantes en $d$ sont $(u_d)_d \circ m = m$ et $m'$. Ainsi $m = m'$. Comme $L(d)$ et $L(d')$ sont deux limites du même diagramme, la flèche $L(f)$, compatible aux cônes, est l'isomorphisme canonique entre elles.

Les colimites se traitent de façon duale. Pour l'objet nul, un objet nul de $\mathcal A^{\mathcal C}$ est la limite et la colimite du diagramme vide, et sa valeur en un objet $d$ est donc à la fois finale et initiale dans $\mathcal A$. Pour la structure abélienne, $\mathcal A$ a les limites et colimites finies par ce qui précède ; la flèche canonique de la coimage vers l'image d'une flèche $f$ de $\mathcal A$ devient, par $\mathrm{const}$, celle de $\mathrm{const}\, f$, qui est un isomorphisme puisque $\mathcal A^{\mathcal C}$ est abélienne. Enfin $\mathrm{const}$ reflète les isomorphismes : l'inverse de $\mathrm{const}\, f$ a pour composante, en n'importe quel objet de $\mathcal C$, un inverse de $f$.
\end{proof}

\noindent\emph{Vérification en Lean.} Le lemme est l'instance \texttt{createsLimitsOfSize\_const}, avec son dual \texttt{createsColimitsOfSize\_const} ; ses étapes sont \texttt{u\_app\_self} (les composantes diagonales sont des identités), \texttt{isLimitEvalCone} ($L(d)$ est une limite) et \texttt{isIso\_pt\_map} ($L$ envoie toute flèche sur un isomorphisme). Les conséquences sont \texttt{hasLimitsOfShape\_of\_functorCategory}, \texttt{hasFiniteLimits\_of\_functorCategory} et leurs duales, \texttt{hasZeroObject\_of\_functorCategory} et \texttt{abelianOfFunctorCategory}, dans \href{\leanblobmain/Redaction469.lean}{\texttt{Redaction469.\allowbreak{}lean}}. Mathlib avait déjà le fait que $\mathrm{const}$ préserve les limites qui existent dans $\mathcal A$ (\texttt{preservesLimits\_const}), mais pas la création.

Ne sont pas démontrés ici : la structure préadditive de $\mathcal A$, que le corollaire abélien suppose donnée au lieu de la tirer de $\mathcal A^{\mathcal C}$ (addition des flèches par les biproduits et l'objet nul, puis opposés) ; le résultat est donc partiel sur ce point. Le théorème~4 lui-même, dans le sens direct, est \texttt{functorCategoryAbelian} de mathlib.

\litterature{le lemme des errata est tenu pour une piste candidate. La recherche consignée se limite à mathlib, qui n'a que la préservation des limites par $\mathrm{const}$.}
